\section*{अध्याय \ref{ch:b1:s-block} --- हाइड्रोजन और s-ब्लॉक}

\begin{solution}{exo:b1:s-block:1}
\ce{KH} और \ce{CaH2}: लवणीय; \ce{SiH4} और \ce{H2S}: सहसंयोजी;
\ce{PdH_{0.6}}: धात्विक। \ce{CaH2 + 2H2O -> Ca(OH)2 + 2H2}।
\end{solution}

\begin{solution}{exo:b1:s-block:2}
\ce{2Li + 2H2O -> 2LiOH + H2}, और Na तथा K के साथ भी यही;
\ce{Mg + 2HCl -> MgCl2 + H2}।
\end{solution}

\begin{solution}{exo:b1:s-block:3}
क्षारकीय: \ce{Na2O}, \ce{MgO} (\ce{MgO + 2HCl -> MgCl2 + H2O})। अम्लीय:
\ce{SiO2}, \ce{SO3}, \ce{CO2} (\ce{SO3 + H2O -> H2SO4})। उभयधर्मी:
\ce{Al2O3}, \ce{ZnO} (\ce{ZnO + 2OH- + H2O -> Zn(OH)4^2-})।
\end{solution}

\begin{solution}{exo:b1:s-block:4}
pH 14 पर जल का युग्म \ce{2H2O + 2e- -> H2 + 2OH-} है, $E = -0.83$~V;
$\log K = 2(-0.83 + 2.71)/0.059 = 64$।
\end{solution}

\begin{solution}{exo:b1:s-block:5}
ज्वाला में संघट्ट किसी इलेक्ट्रॉन को उत्तेजित स्तर पर उठा देते हैं; \omterm{def:b1:quantum-numbers:energy-level}{मूल अवस्था} में
लौटते समय परमाणु अंतराल के बराबर ऊर्जा का एक फ़ोटॉन उत्सर्जित करता है, जो
तत्त्व के लिए विशिष्ट है। $E = hc/\lambda = 6.626 \times 10^{-34} \times
3.00 \times 10^8/(589 \times 10^{-9}) = \qty{3.375e-19}{J} = \qty{2.11}{eV}$।
% ledger: const:h, const:c, const:eV
\end{solution}

\begin{solution}{exo:b1:s-block:6}
$10^6/106.0 = \qty{9.43e3}{mol}$ सोडा ऐश: \qty{9.43e3}{mol}
चूना पत्थर, \qty{0.944}{t}, और $2 \times 9.43 \times 10^3/0.90 =
\qty{2.10e4}{mol}$ लवण, \qty{1.23}{t}।
\end{solution}

\begin{solution}{exo:b1:s-block:7}
$[\ce{Mg^2+}] = 1.3/24.3 = \qty{0.053}{mol/L}$; $\mathrm{pH} = 14.00 -
\frac12(11.25 + \log 0.053) = 9.0$। \qty{1.0}{m^3} में मैग्नीशियम के \qty{53.5}{mol}
को \ce{Ca(OH)2} के \qty{53.5}{mol} चाहिए: \qty{4.0}{kg}।
\end{solution}

\begin{solution}{exo:b1:s-block:8}
साम्य स्थिरांक (ऊष्मागतिकी) दोनों के लिए विशाल है; दर
(गतिकी) का निर्णय $E^\circ$ से नहीं होता। लिथियम का गलनांक
अधिक है, वह ठोस रहता है और कम विलेय उत्पाद की परत से ढक जाता है,
और पोटैशियम की तुलना में अपना इलेक्ट्रॉन कम आसानी से खोता है (उच्चतर आयनन
ऊर्जा): वह स्थिर गति से अभिक्रिया करता है, पोटैशियम प्रचंडता से।
\end{solution}

\begin{solution}{exo:b1:s-block:9}
$10^6/100.1 = \qty{9.99e3}{mol}$: CaO $9.99 \times 10^3 \times 56.1 =
\qty{0.560}{t}$; \ce{CO2} \qty{9.99e3}{mol}, $V = nRT/p = 9.99 \times 10^3
\times 8.314 \times 298.15/10^5 = \qty{248}{m^3}$।
\end{solution}

\begin{solution}{exo:b1:s-block:10}
$2.9 \times 10^8/6.9 = \qty{4.2e7}{kmol}$ Li (द्रव्यमान kg में), $\times 73.8/2 =
\qty{1.55e9}{kg}$: लगभग 16 लाख टन लिथियम कार्बोनेट।
$2.9 \times 10^8/8 = \num{3.6e7}$ बैटरियाँ।
\end{solution}

\begin{solution}{exo:b1:s-block:11}
20, 80, \qty{100}{\celsius} पर 1.31, 0.84, 0.71 प्रतिशत: \omterm{def:b1:precipitation:solubility}{विलेयता}
घटती है। गरम अवस्था में कम कार्बोनेट विलयन में रहता है: अधिक पुनःप्राप्त होता है।
\end{solution}

\begin{solution}{exo:b1:s-block:12}
आयन जो लवण बना सकते हैं उनमें \ce{NaHCO3} सबसे कम विलेय है, और
विलयन सबसे पहले उसी से संतृप्त होता है। अमोनिया विलयन को इतना क्षारकीय बना देती है
कि \ce{CO2}, \ce{HCO3-} में बदल जाए, और साथ ही वह पुनःप्राप्त की जा सकती है
(\ce{NH4Cl} के रूप में, जिसे चूना पुनर्जनित करता है); सोडियम हाइड्रॉक्साइड खप जाता और
उत्पाद से अधिक महँगा पड़ता।
\end{solution}

\begin{solution}{pb:b1:s-block:1}
\textbf{1.} Li \qty{6.0}{g/L}, Mg \qty{2.4}{g/L}।
\textbf{2.} चार घन मीटर लवण-जल एक देते हैं: तीन घन मीटर
जल वाष्पित होता है।
\textbf{3.} धूप और शुष्क जलवायु ऊर्जा मुफ़्त देते हैं।
\textbf{4.} सोडियम क्लोराइड, जो कहीं सबसे प्रचुर लवण है, सबसे पहले
संतृप्त होता है।
\textbf{5.} Li $6000/6.9 = \qty{870}{mol}$; Mg $2400/24.3 = \qty{98.8}{mol}$।
\textbf{6.} मैग्नीशियम कार्बोनेट के साथ अवक्षेपित होकर
उत्पाद को संदूषित कर देगा।
\textbf{7.} \ce{Mg^2+ + Ca(OH)2 -> Mg(OH)2 + Ca^2+}।
\textbf{8.} $98.8 \times 74.1 = \qty{7.32}{kg}$।
\textbf{9.} बचा हुआ $[\ce{Mg^2+}] = \num{9.88e-5}$ mol/L: $\mathrm{pH} = 14.00 -
\frac12(11.25 - 4.01) = 10.38$।
\textbf{10.} लिथियम हाइड्रॉक्साइड विलेय है: इस pH पर कोई ठोस नहीं बनता।
\textbf{11.} \ce{Ca^2+ + CO3^2- -> CaCO3}।
\textbf{12.} $98.8 \times 106.0 = \qty{10.5}{kg}$।
\textbf{13.} \ce{2Li+ + CO3^2- -> Li2CO3}।
\textbf{14.} $435 \times 106.0 = \qty{46.1}{kg}$।
\textbf{15.} $435 \times 73.8 = \qty{32.1}{kg}$।
\textbf{16.} लिथियम कार्बोनेट गरम में कम विलेय है।
\textbf{17.} $0.0084 \times 1000 = \qty{8.4}{kg}$।
\textbf{18.} गरम जल की थोड़ी मात्राओं से, जिसमें वह सबसे कम
विलेय है।
\textbf{19.} $32.1 - 8.4 = \qty{23.7}{kg}$।
\textbf{20.} $23.7/32.1 = \qty{74}{\percent}$।
\textbf{21.} उसमें अब भी लिथियम है: उसे तालाबों में लौटा दिया जाता है।
\textbf{22.} $2.9 \times 10^8/6.0 = \num{4.8e7}$ घन मीटर।
\textbf{23.} प्रति घन मीटर \textbf{\qty{24}{kg}} (23.7) लिथियम कार्बोनेट।
\end{solution}
